\Alg

\begin{enumerate}
    \item Отсортируем все точки по абсциссе.
    \item Разделим медианой по абсциссе точки на две группы.
    \item Найдем в каждой рекурсивно ответ.
    \item Сольем две половины и найдем оптимальный ответ.
\end{enumerate}

Пусть $d\,=\,min(d_1,\,d_2)$, где $d_1$ — ответ от левой половины, а $d_2$ — от правой.
Ответ может быть либо равен d, либо надо рассмотреть пары точек из обеих частей.
Заметим, что если мы разбили по $x_0$ пополам, то надо рассматривать только точки с абсциссой в пределах $[x_0\,-\,d,\,x_0\,+\,d]$

Пусть $P_i$ попадает в полосу, обозначим за $C(P_i)$ прямоугольник по
абсциссе в $[x_0\,-\,d,\,x_0\,+\,d]$, а по ординате $[P_i.y,\,P_i.y\,+\,d]$.\\
Разобьем $C(P_i)$ на 8 квадратов $\frac{d}{2} \times \frac{d}{2}$. Не может быть двух точек в одном из квадратов, так как иначе расстояние на одной из половин было бы меньше $d$.

\drawsomesmall{pix/search_a_pair_of_nearby_points.png}

Откуда для каждой точки в полосе $O(1)$ точек кандидатов.\\
Отсортируем точки по ординате и будем идти двумя указателями по
тем что в полосе, ища минимальное расстояние.

\subsection*{Асимптотика}

Заметим, что если время слияния двух мнодеств по n точек равно $Merge(n)$, то итоговое время работы: $T(n)\,=\,2T(n/2)\,+\,Merge(n)$.\\
По мастер-теореме $T(n)\,=\,O(Merge(n)\, \cdot \,log\,n)$.\\
В слиянии сначала происходит сортировка, а потом два указателя, по
$O(1)$ времени на каждое положение, а значит $Merge(n)\,=\,O(n\,log\,n)$.\\
То есть время работы равно $O(n\,log^2\,n)$.

\Note Если алгоритм переупорядочивает набор точек, уже отсортированных
по y, то вместо того, чтобы заново сортировать набор на стадии слияния применим Merge за линейное время!
Тогда $Merge(n)\,=\,O(n)$ и итоговое время работы равно $O(n\,log\,n)$.

\pagebreak